\section{Worked Examples and Comparative Case Studies}

This section is meant to consolidate the whole chapter.
Up to this point, convolution, recurrence, gating, attention, and transformers have been introduced as different architectural responses to different kinds of structure.
The most useful final step is to see these ideas operating on small, concrete examples.
A worked example forces us to ask the right question: what structural assumption is the model making, what operator is it applying, and why is that operator well matched --- or badly matched --- to the task?

\medskip

The examples below are deliberately small.
Their purpose is not to impress by scale, but to reveal the logic of the architecture.
In each case, one should pay attention not only to the numerical computation, but also to the inductive bias that computation expresses.

\subsection*{A mini CNN example: detecting a vertical edge}

Consider the following toy grayscale image:
\[
X=
\begin{bmatrix}
1 & 1 & 0 & 0\\
1 & 1 & 0 & 0\\
1 & 1 & 0 & 0\\
1 & 1 & 0 & 0
\end{bmatrix}.
\]
The image is bright on the left and dark on the right, so it contains a vertical transition near the middle columns.
A natural convolutional filter for this kind of pattern is
\[
K=
\begin{bmatrix}
1 & -1\\
1 & -1
\end{bmatrix}.
\]
This filter compares the left and right parts of a small patch.
If both sides of the patch are similar, the output is small.
If the left side is bright and the right side is dark, the output is large and positive.

Using stride $1$ and no padding, the convolution output has size $3 \times 3$.
Each entry is obtained by placing the filter on a $2 \times 2$ patch and taking the inner product.
For the first row of patches we obtain
\[
\langle K, X_{1:2,1:2}\rangle
=
\left\langle
\begin{bmatrix}1 & -1\\ 1 & -1\end{bmatrix},
\begin{bmatrix}1 & 1\\ 1 & 1\end{bmatrix}
\right\rangle
=1-1+1-1=0,
\]
\[
\langle K, X_{1:2,2:3}\rangle
=
\left\langle
\begin{bmatrix}1 & -1\\ 1 & -1\end{bmatrix},
\begin{bmatrix}1 & 0\\ 1 & 0\end{bmatrix}
\right\rangle
=1+0+1+0=2,
\]
\[
\langle K, X_{1:2,3:4}\rangle
=
\left\langle
\begin{bmatrix}1 & -1\\ 1 & -1\end{bmatrix},
\begin{bmatrix}0 & 0\\ 0 & 0\end{bmatrix}
\right\rangle
=0.
\]
By the same reasoning for the remaining rows, the full feature map is
\[
X * K=
\begin{bmatrix}
0 & 2 & 0\\
0 & 2 & 0\\
0 & 2 & 0
\end{bmatrix}.
\]
After a ReLU nonlinearity, nothing changes because all entries are already nonnegative.
A max-pooling or global-max step would then return the value $2$, signaling that the image contains the desired vertical edge pattern somewhere in the field of view.

\begin{example}[What this example shows]
The important point is not the specific number $2$.
The important point is that the \emph{same} filter was applied at every location.
The model did not need one parameter set for the top edge, another for the middle edge, and another for the bottom edge.
That is the inductive bias of convolution: local pattern detectors are shared across space.
If the same kind of edge appears in another location, the same filter can detect it there too.
\end{example}

A dense fully connected layer could also learn to detect this image pattern, but it would not build translation sensitivity and parameter sharing into the operator itself.
The CNN does so directly.
This is why convolution is especially natural for images: the operator already encodes locality and repeated spatial motifs.

\subsection*{Recurrence versus attention on a toy sequence memory task}

Now consider a toy sequence problem.
A sequence consists of symbols from the set $\{A,B\}$, and the target at the final position is: \emph{repeat the first symbol of the sequence}.
For example,
\[
(A,B,B,A,\, ?)
\longmapsto A.
\]
This is a deliberately simple task, but it separates two architectural ideas very clearly.
The model must preserve or recover information from the beginning of the sequence when it reaches the end.

\paragraph{Recurrent solution.}
A recurrent model computes
\[
h_t = \phi(W_h h_{t-1} + W_x x_t),
\]
where $x_t$ is the embedding of the current token and $h_t$ is the hidden state.
To solve the task, the network must somehow store the information contained in $x_1$ inside the evolving state and prevent later updates from destroying it.
For a short sequence this is possible.
For a long sequence it becomes harder, because the hidden state must continuously compress the past into a single vector while still processing new inputs.
The required information travels through the chain
\[
x_1 \to h_1 \to h_2 \to h_3 \to \cdots \to h_T.
\]
Thus the relevant interaction path from the first token to the last output has length on the order of $T$.

\paragraph{Attention-based solution.}
A decoder-style or self-attention model can treat the same task differently.
At the final position, it may compute
\[
y_T = \sum_{j=1}^{T} a_{Tj} v_j,
\qquad
\sum_{j=1}^{T} a_{Tj}=1,
\]
where $v_j$ is the value vector for token $j$.
If the model learns that the final query should focus on the first position, then one possible attention pattern is
\[
(a_{T1},a_{T2},a_{T3},a_{T4},a_{T5})
=(0.92,0.03,0.02,0.02,0.01).
\]
Then
\[
y_T \approx 0.92 v_1 + 0.03 v_2 + 0.02 v_3 + 0.02 v_4 + 0.01 v_5,
\]
so the output is dominated by the representation of the first symbol.
The relevant information need not survive repeated state updates; it can be retrieved directly.

\begin{example}[Same task, different operator]
For this toy task, the recurrent architecture behaves like a running memory state, while self-attention behaves like a content-based retrieval rule.
The RNN says: ``compress the past into a state and carry it forward.''
Attention says: ``when the final query is formed, look back and pull the needed representation directly.''
Neither statement is universally better in every setting, but they represent sharply different inductive biases.
\end{example}

This example also clarifies why gated recurrence was historically important.
An LSTM or GRU modifies the recurrent update so that important information can be preserved more selectively, partially narrowing the gap between plain recurrence and attention-based access.
Still, the underlying logic remains different: recurrence transports information through state evolution, while attention creates direct interaction edges.

\subsection*{A toy machine-translation alignment example}

Machine translation is a particularly instructive case because it requires both representation and alignment.
Suppose the source sentence is
\[
x=(\text{le},\text{chat},\text{dort})
\]
and the target sentence is
\[
y=(\text{the},\text{cat},\text{sleeps}).
\]
Let the encoder produce source representations
\[
H=
\begin{bmatrix}
h_1\\ h_2\\ h_3
\end{bmatrix}
\in \mathbb{R}^{3 \times d},
\]
corresponding to the three source words.
At target step $i$, the decoder forms a query $q_i$ and computes cross-attention weights
\[
\alpha_i = \mathrm{softmax}\!\left(\frac{q_i K^{\top}}{\sqrt{d_k}}\right),
\qquad K=HW_K,
\qquad V=HW_V.
\]
Equivalently, for all target positions at once,
\[
C = \mathrm{softmax}_{\mathrm{row}}\!\left(\frac{QK^{\top}}{\sqrt{d_k}}\right)V,
\]
where the rows of $Q$ come from target-side decoder states.

To make this concrete, suppose the decoder produces the following score matrix before softmax:
\[
S=
\begin{bmatrix}
3 & 1 & 0\\
0 & 4 & 0\\
0 & 0 & 4
\end{bmatrix}.
\]
Applying row-wise softmax gives approximately
\[
A=\mathrm{softmax}_{\mathrm{row}}(S)
\approx
\begin{bmatrix}
0.84 & 0.11 & 0.05\\
0.02 & 0.96 & 0.02\\
0.02 & 0.02 & 0.96
\end{bmatrix}.
\]
The interpretation is immediate:
\begin{itemize}
    \item when producing \emph{the}, the decoder attends mostly to \emph{le},
    \item when producing \emph{cat}, it attends overwhelmingly to \emph{chat},
    \item when producing \emph{sleeps}, it attends overwhelmingly to \emph{dort}.
\end{itemize}
The corresponding context vectors are
\[
c_1 = 0.84v_1+0.11v_2+0.05v_3,
\]
\[
c_2 = 0.02v_1+0.96v_2+0.02v_3,
\]
\[
c_3 = 0.02v_1+0.02v_2+0.96v_3.
\]
These are not hard dictionary lookups.
They are soft weighted combinations of encoded source representations.
That softness matters because translation is rarely a one-word-at-a-time substitution rule.
Sometimes multiple source words contribute jointly, and sometimes word order differs between languages.

\begin{example}[Why attention solved a real bottleneck]
In a classical encoder--decoder model without attention, the entire source sentence had to be compressed into a single fixed-size vector before decoding began.
For longer or more complex sentences, this creates an information bottleneck.
Cross-attention removes much of that bottleneck by allowing each target token to consult the full source representation sequence directly.
Thus attention is not just a convenience for visualization; it solves a precise architectural problem.
\end{example}

This machine-translation example also shows why alignment matrices became such powerful pedagogical objects.
They let us inspect, at least approximately, which source positions influenced which target outputs.
That makes the operator visible in a way that a hidden recurrent state often is not.

\subsection*{A time-series forecasting example}

Consider a one-step-ahead forecasting problem for a short noisy time series:
\[
(20,22,21,23).
\]
Suppose we believe the underlying signal evolves through a slowly changing latent level and that short-term noise should be smoothed.
A very simple recurrent state model for this idea is
\[
h_t = 0.7 h_{t-1} + 0.3 x_t,
\qquad \hat{x}_{t+1}=h_t,
\]
with initial state $h_1=x_1=20$.
Then
\[
h_2 = 0.7\cdot 20 + 0.3\cdot 22 = 20.6,
\]
\[
h_3 = 0.7\cdot 20.6 + 0.3\cdot 21 = 20.72,
\]
\[
h_4 = 0.7\cdot 20.72 + 0.3\cdot 23 \approx 21.40.
\]
So the one-step-ahead forecast is
\[
\hat{x}_5 = h_4 \approx 21.40.
\]

This forecast does not simply copy the last observation $23$.
Instead it forms a smoothed state that reflects the recent history.
That is precisely the inductive bias of recurrent or state-space style forecasting: observations are treated as partial evidence about an evolving hidden state.

Now compare this with a causal convolutional rule such as
\[
\hat{x}_{t+1} = 0.5x_t + 0.5x_{t-1}.
\]
Applied at $t=4$, this gives
\[
\hat{x}_5 = 0.5\cdot 23 + 0.5\cdot 21 = 22.
\]
This rule is also reasonable, but it encodes a different assumption.
It says that the next value depends mainly on a local window of recent observations, not on a hidden state that is updated over time.
An attention-based forecaster would express yet another assumption: the next value may depend on selected past times, perhaps even distant ones, if their patterns are relevant.

\begin{example}[What the forecasting example teaches]
There is no unique ``time-series architecture.''
If the series is governed by a slowly evolving latent state, recurrence is very natural.
If short local motifs dominate, causal convolution may be simpler and more effective.
If irregular long-range interactions matter --- for example, if today should be compared with the same day in earlier cycles --- attention can be very attractive.
The architecture should follow the dependence structure one expects in the data.
\end{example}

\subsection*{Final comparison across the chapter}

The worked examples point toward one central lesson.
Architectural names are secondary.
What matters is the structural claim each architecture makes about the data and the operator by which it uses that claim.
Table~\ref{tab:arch-compare} summarizes the chapter in that language.

\begin{table}[ht]
\centering
\scriptsize
\setlength{\tabcolsep}{4pt}
\begin{tabular}{|p{0.14\textwidth}|p{0.16\textwidth}|p{0.16\textwidth}|p{0.14\textwidth}|p{0.14\textwidth}|p{0.11\textwidth}|}
\hline
\textbf{Architecture} & \textbf{Assumption} & \textbf{Operator} & \textbf{Strength} & \textbf{Weakness} & \textbf{Natural domain} \\
\hline
CNN & Local spatial structure and repeated patterns across location & Shared local convolution filters, often followed by pooling & Parameter sharing; strong locality prior; translation-friendly features & Less natural for arbitrary long-range dependence without added mechanisms & Images, spatial signals \\
\hline
RNN & Sequence is processed through evolving state & Recurrent update $h_t=\phi(W_hh_{t-1}+W_xx_t)$ & Natural state evolution; compact online processing & Long-range information must survive many updates; sequential computation & Time series, streaming sequences \\
\hline
LSTM / GRU & Important state information should be selectively preserved or overwritten & Gated recurrent state transitions & Better memory control than plain RNNs; robust on many sequential tasks & Still fundamentally sequential; state bottleneck remains & Language, speech, temporal data \\
\hline
Attention / encoder--decoder attention & Relevance should be computed content-dependently between positions or streams & Weighted aggregation from learned similarity scores & Explicit alignment; direct access to relevant source information & Dense interaction can be expensive; needs careful representation design & Translation, retrieval, cross-modal fusion \\
\hline
Transformer & Rich sequence behavior can be built from repeated attention blocks plus MLPs, residuals, and normalization & Multi-head self-attention and feedforward blocks & Parallelizable global interaction; flexible across many modalities & Dense attention scales quadratically in context length; positional structure must be supplied & Language, vision patches, multimodal token streams \\
\hline
\end{tabular}
\caption{A comparative summary of the main architectures in this chapter. The key comparison is not historical order, but the relation between structural assumption and operator.}
\label{tab:arch-compare}
\end{table}

The table should not be read as a ranking.
It is a map.
Each row describes a way of embedding prior assumptions into a learnable system.
The central unifying message of the chapter is therefore the following: architecture is the mechanism by which inductive bias becomes operational.

\paragraph{What to remember.}
Convolution, recurrence, gating, attention, and transformers should be compared through the structures they assume and the operators they apply.
A CNN uses shared local filters and is naturally suited to repeated spatial patterns.
A recurrent model uses an evolving state and is naturally suited to temporal evolution.
Attention uses content-dependent weighted aggregation and is naturally suited to alignment and selective retrieval.
A transformer turns attention into a full repeated architectural block that supports flexible global interaction.
The best architecture is not the most fashionable one, but the one whose operator matches the dependence structure of the task.

\paragraph{Common confusion.}
Students often ask which architecture is ``best.''
That question is usually too vague.
A better question is: what does the task require the model to preserve, compare, or aggregate?
It is also easy to confuse representational power in principle with suitability in practice.
A dense network can represent many of the same functions as a CNN or transformer, but it may do so with a much weaker built-in prior and much less efficient learning.
Finally, one should not confuse a worked toy example with a full empirical benchmark.
The role of a toy example is to reveal the structural logic of the operator.

\paragraph{Exercises.}
\begin{enumerate}
    \item In the CNN example, replace the filter
    \[
    \begin{bmatrix}1 & -1\\ 1 & -1\end{bmatrix}
    \]
    by
    \[
    \begin{bmatrix}1 & 1\\ -1 & -1\end{bmatrix}.
    \]
    Compute the new feature map and explain what pattern this second filter is designed to detect.
    \item For the toy memory task, explain from first principles why the information path from the first token to the last output has length on the order of $T$ in a recurrent model but length $1$ in a self-attention layer.
    \item In the translation example, verify numerically that each row of the alignment matrix $A$ sums to one.
    Then explain why this implies that each context vector is a convex combination of value vectors.
    \item Compare the recurrent forecasting rule
    \[
    h_t=0.7h_{t-1}+0.3x_t
    \]
    with the causal convolutional rule
    \[
    \hat{x}_{t+1}=0.5x_t+0.5x_{t-1}.
    \]
    What hidden assumption about temporal structure is built into each method?
    \item Give one task for which a CNN is likely to be more natural than a transformer, and one task for which a transformer is likely to be more natural than a plain RNN.
    Justify both answers in terms of inductive bias rather than popularity.
    \item Choose any two rows from Table~\ref{tab:arch-compare} and write a short proof-style argument explaining why their operators impose genuinely different constraints on how information can move through the network.
\end{enumerate}

\paragraph{Notes and references.}
The convolutional viewpoint developed in this chapter traces back to the early CNN literature associated with LeCun and collaborators.
The recurrent and gated examples are closely connected to the RNN, LSTM, and GRU traditions developed by Elman, Hochreiter and Schmidhuber, and Cho and collaborators.
The translation alignment example reflects the conceptual importance of attention in neural machine translation, especially in the work of Bahdanau, Cho, and Bengio.
The transformer synthesis inherits the operator viewpoint of Vaswani and collaborators and later general-purpose sequence modeling work.
Pedagogically, these examples are best read not as historical trophies, but as comparative demonstrations of how architectural assumptions become computational operators.
